% TN0484 / WP0185: abstract included by tn0484_storyline.tex and mirrored in the TN manuscript.
% Date: 2026-10-03. v0.14.5 implements K's approved weak-field effective-response framing.
% The JNE build keeps a structured abstract with the same scope. Section and figure order stay fixed.
Temporal-interference (TI) stimulation applies two high-frequency currents whose
superposition contains no field component at their low difference frequency. Population
models generate that component through effective quadratic mixing and amplify it by
recurrent resonance. For a smooth weak-input response with a nonzero second-order term,
cellular nonlinearities contribute to a beat proportional to the product of the carrier
amplitudes. A Jansen--Rit firing-rate sigmoid realizes this mixing through its curvature.
An exact quadratic integrate-and-fire mean field provides a derived realization: its
static rate curve emerges when membrane relaxation is fast relative to the input, while
synapses remain dynamic. Outside that limit the mixing coefficient is frequency
dependent. Fast cellular access and slow circuit resonance can therefore coexist, with
membrane attenuation counted once. Near a Hopf bifurcation, the network amplifies the
generated beat as $1/\gamma$ within linear response, where $\gamma$ is the modal damping.
Coupling sweeps separate this gain from leading-order beat generation. Beat-locked rate
modulation grows faster than the mean-rate change in the mean field and in a spiking
network of $8000$ neurons per population. In this weak-input regime, the carrier-product
law makes an envelope-only dose map conditional on the transduction mechanism. TI and
direct low-frequency stimulation share the amplifier when effective input coupling,
readout, and operating state are matched. The analysis concerns weak modulation of
ongoing activity; biological access and the absolute response at human field strengths
remain unresolved.
